\scanpage{6}
Soit $G_0$ l'image de $\unsure{\mathscr G_n|G_n|\mathfrak G_n}$ dans $\operatorname{Aut}(E,\deleted{\unsure{f},}V)$, ou mieux dans $\operatorname{Aut}(E,f,\pi,\pi')$
\diagram{\begin{tikzcd}[ampersand replacement=\&,column sep=large]
\deleted{G_0}\arrow[r,draw=none,"\deleted{\longrightarrow}" description]\&\operatorname{Aut}(X,C)\deleted{\unsure{Q_X}}\arrow[r]\&\operatorname{Aut}(C,\deleted{\unsure{Q_C}})\\
G_0\arrow[r,hook]\&\operatorname{Aut}(E,V,f)\arrow[u]\arrow[r]\&\operatorname{Aut}(V,\deleted{f_V})\arrow[u]
\end{tikzcd}}
On cherche \deleted{les} $\inserted{3}$ noyaux de
\diagram{\begin{tikzcd}[ampersand replacement=\&,column sep=small,row sep=small]
\&[-1em] \operatorname{Aut}(X,\deleted{\unsure{C}})\arrow[dr]\&\\
G_0\arrow[r]\&\left\{\vphantom{\dfrac{A}{A}}\right.\&\operatorname{Aut}(C)\\
\&\operatorname{Aut}(V)\arrow[ur]\&
\end{tikzcd}}

a) $\operatorname{Ker}(\operatorname{Aut}(E,V,f)\deleted{\unsure{\to}}\longrightarrow\operatorname{Aut}(X))$ est formé des homothéties qui conservent $f$, donc c'est $\mu_2$. En fait, $G_0$ conserve mieux que $(E,f,V)$, il conserve $(E,f,\pi,\pi')$ \inserted{\unsure{pour n'importe $f$|peu importe $f$}}, donc si c'est une homothétie c'est l'identité : $G_0$ opère \uline{fidèlem$^{\mathrm t}$} sur $X$.

b) $\operatorname{Ker}(\operatorname{Aut}(E,f,\pi,\pi')\longrightarrow\operatorname{Aut}(V))$ est formé des morphismes \inserted{$u$} de la forme \deleted{\unsure{$u$}} $\mathrm{id}+\deleted{\unsure{\lambda}}\pi'\otimes a$ $(a\in\Gamma E)$, qui conservent $\pi,\pi',f$. Or $u$ conserve $\pi$ sss $\deleted{\unsure{\lambda}}\rho a=0$ (donc $a=0$ si $\rho$ inv., plus gén. si $\rho$ est \uline{régulier}) (NB $\langle\pi,\pi'\rangle=\rho$), et il conserve $\pi'$ sss :
\[
\langle a,\pi'\rangle\deleted{\unsure{\lambda}}\ \pi'=0\qquad\text{i.e.}\quad\deleted{\unsure{\lambda}}\langle a,\pi'\rangle=0\qquad\text{i.e.}
\]
$\boxed{a\in V.}$ Conserver $f$ signifie
\[
f(a)\pi'+\psi(a)=0\tag*{(\unsure{*|8})}
\]
d'où \deleted{\unsure{nous}} en appliquant $\psi'$
\[
f(a)\underbrace{\psi'(\pi')}_{-\sigma'\pi}+\underbrace{\psi'\psi(a)}_{\beta a}=0\qquad\text{i.e.}\quad\beta a=\deleted{\unsure{h}}\sigma'f(a)\pi
\]

% archive sheet number: 5

